\fancyhead[L]{\small ACHROMATIC DOUBLET OBJECTIVES.}
\SourceLeaf{510}\pdfbookmark[0]{On Achromatic Doublet Objectives}{achromatische_doppelobjective}
\begin{center}
{\Large ON ACHROMATIC DOUBLET OBJECTIVES\par}
\vspace{6mm}
{\large WITH PARTICULAR REGARD TO\par}
\vspace{4mm}
{\large THE COMPLETE ELIMINATION OF COLOUR DISPERSION.\par}
\vspace{7mm}\hrule\vspace{2mm}
Zeitschrift für Astronomie und verwandte Wissenschaften\par
edited by B.\ von \textsc{Lindenau} and \textsc{Bohnenberger}. Vol.\ IV. No.\ XXX. December 1817.
\vspace{2mm}\hrule\vspace{5mm}
\end{center}
Professor \textsc{Bohnenberger}'s fine paper on achromatic objectives in the first volume of this journal has the merit of being the first to raise a point important to this theory. It has prompted me to take up some earlier investigations again and develop them further. I shall report the results here.

Until now, in doublet objectives, it was considered sufficient to eliminate colour dispersion for rays infinitely close to the axis and spherical aberration for rays of intermediate refrangibility. This can still leave some colour dispersion for the marginal rays. With this arrangement, calculating an achromatic objective is an indeterminate problem. That is, for any crown-glass lens with a positive focal length, whatever the ratio of the radii of its surfaces, a flint-glass lens can be calculated that, when combined with the crown-glass lens, gives an objective that is achromatic in the sense just described. As far as I know, all opticians have so far assumed both surfaces of the crown-glass lens to be convex. Theorists, however, have proposed very different values for the ratio of the two radii, depending on the principle they started from. If, with \textsc{Euler}, one wishes to make the aberration due to the shape of the crown-glass lens as small as possible, the ratio of the radii must be about
\SourceLeaf{511}1 to 7. The radii must be equal if one wants the smallest possible curvatures, as \textsc{Klügel} does in his analytical dioptrics. If the refractions themselves are to be as small as possible, as the same writer intends in a later paper, they must be equal, and the radii must be in a ratio close to 1 to 3. It does not seem that all these different proposals are sufficiently justified. \textsc{Klügel}'s main concern was spherical aberration, which, as is well known, cannot be eliminated with mathematical exactness for \textit{all} rays. In \textsc{Euler}'s treatment of these calculations, this aberration is really eliminated only for the rays nearest the axis. The remaining aberration is very nearly proportional to the fourth power of the distance from the axis, and is therefore most noticeable for the marginal rays. Or if one follows \textsc{Klügel} and does the calculation so that the aberration vanishes for the marginal rays, it reappears for the intermediate rays, most noticeably for those whose distance from the axis is nearly $\tfrac{7}{10}$ of the aperture radius. In choosing the ratio of the first two radii, \textsc{Klügel}'s aim was to make this unavoidable residual aberration due to shape as harmless as possible. It is not clear enough, however, that the chosen ratio really achieves this aim best, or that the aim is important enough to give it priority as the sole basis for determining the ratio. For if such an objective has \textit{other} imperfections substantially greater than those due to the aberration arising from shape that cannot be eliminated entirely, it is evidently more important to take the former into account than the latter.

For this reason it will be better to use the freedom we have in determining the ratio of the first two radii to reduce or eliminate colour dispersion for the marginal rays. Professor \textsc{Bohnenberger} has in fact shown by calculation that in this respect the ratio of 2 to 3 is preferable to that of 1 to 3. The first produces considerably less colour dispersion for the marginal rays without making the remaining spherical aberration appreciable. Nevertheless, Professor \textsc{Bohnenberger}'s arrangement also leaves some colour dispersion for the marginal rays,
\SourceLeaf{512}which it would be very desirable to reduce further or eliminate altogether. According to his statement on p.\ 392, Professor \textsc{Bohnenberger}'s experiments for this purpose were unsuccessful and might seem to warrant the supposition that this was impossible. This prompted me to investigate the matter separately, and I found the opposite, which seems to me quite remarkable.

Complete elimination of colour dispersion for the marginal rays and those nearest the axis is indeed possible. More precisely, an objective can be calculated that brings all rays of two specified colours to one and the same point, both those that strike it at a specified distance from the axis and those that strike it infinitely close to it, parallel to the axis as is always assumed here. This gives the objective a shape quite different from those used exclusively up to now. Both lenses are convex-concave, with their convex surfaces facing the object. Although this results in greater refractions than in other arrangements, the unavoidable aberration due to shape that remains is still very slight. As a result, \textit{all} rays that strike the objective parallel to the axis are brought together more perfectly than in any other arrangement. It would therefore be well worth having skilled lens makers try this new shape. There may still be \textit{practical} difficulties at present. One of them will be that the pieces of glass from which the lenses are to be ground will have to be thicker. But with the continual improvement of the technical side of dioptrics, one can hope that such difficulties will be overcome. It will then be up to mathematics to give the ideal shape for bringing the rays together as perfectly as possible.

I give my calculation solely as an example to confirm what I have said. I am not asking lens makers to follow these dimensions exactly. It is essential first to determine the ratios of refraction and dispersion for the kinds of glass from which a perfect objective is to be ground, separately and as accurately as possible, and then to calculate the dimensions of the objective anew accordingly. My calculation used exactly the same figures as
\SourceLeaf{513}\expandafter\def\csname HistoricalNoteText513\endcsname{[Thickness of the first lens $=200$, of the second $=80$ distance between the two lenses $=50$. Values of the refractive ratios for crown glass and flint glass respectively]
\[
\left\{\begin{array}{lll}
1.525976&1.62173&\text{viol. rays}\\
1.515162&1.60177&\text{interm.}\quad\text{—}\\
1.504348&1.58181&\text{red}\quad\text{—}
\end{array}\right\}
\]
}Professor \textsc{Bohnenberger}'s. I also kept the same lens thicknesses and separation\HistoricalNote{513}. But in the new arrangement, the convex surface of the flint-glass lens has a stronger curvature than the concave surface of the crown-glass lens, so the two lenses can be brought closer together (which will also be advantageous in another respect that I will leave aside here). Indeed, if the lens makers have no other objections, the gap can be dispensed with altogether, or the lenses can touch on the axis. Of course, this will entail some modification of the radii of curvature.

I do not intend to develop the mathematical part of this investigation here. I only observe that if the aberration due to shape is treated in \textsc{Euler}'s manner and the thickness and separation of the glass lens are set aside, the problem leads to an equation of degree four with two real roots. The approximate solution obtained from it serves as the basis for an indirect calculation that brings everything into exact agreement. This indication will suffice for mathematicians. One of the real roots of that equation must be discarded, because it entails excessive curvatures of the glass surfaces and would make the incomplete elimination of the aberration due to shape too noticeable.

The result of my calculation is as follows.

If the radii are taken in order as
\[
+3415{,}287;\quad -10133{,}007;\quad +4207{,}421;\quad -2807{,}320
\]
then the red and violet rays, both those that strike infinitely close to the axis and those that strike at a distance of $1083{,}687$, all meet at one point on the axis, whose distance from the last surface is $=28293{,}3$. If the distance from the axis at which the angle of incidence is $18^0\,30'$ is taken as the aperture radius, the diameter of the aperture is very nearly $\tfrac{1}{13}$ of the focal length. To ascer\SourceLeaf{514}tain how large the residual aberration due to shape is for rays between the edge and the axis, I calculated the distances at which the rays meet for an angle of incidence of $13^0$ and found
\[
\begin{array}{ll}
28289{,}3&\text{for the red}\\
28290{,}0&\text{for the violet rays.}
\end{array}
\]

I cannot help adding a remark here on a statement by Professor \textsc{Bohnenberger} in the paper mentioned above. I believe that it is best to eliminate the aberration due to shape exactly for the marginal rays. On p.\ 279 Professor \textsc{Bohnenberger} criticizes this procedure, wrongly in my view. He says that if the aberration were eliminated exactly for the marginal rays, one could enlarge the aperture without impairing clarity up to the point where the aberration again equalled the greatest aberration of the intermediate rays. He therefore favours eliminating the aberration for rays between the edge and the axis instead of for the marginal rays. This would indeed be true if the remaining aberration on either side of the distance at which it is eliminated had \textit{the same sign}, but that is not the case. One might object, with some appearance of plausibility, that the sign of the longitudinal aberration makes no difference at all and that positive and negative aberrations produce one and the same blur in the eye. But that would evidently assume, without saying so, that the eyepiece is always set to show clearly the very image produced by the rays nearest the axis. That assumption cannot be granted. We may call this image the principal image. It coincides with the image produced by the marginal rays if their aberration is eliminated. All the other images then lie, generally speaking, beyond or short of the principal image. Since the eyepiece is always set to make the blur as small as possible, one sees. the principal image least clearly of all, and any enlargement of the aperture also increases the blur. A fuller discussion of this point would take me too far afield here.
